\documentclass[11pt]{article}
\usepackage[letterpaper,margin=0.78in]{geometry}
\usepackage{amsmath,amssymb,mathtools}
\usepackage{booktabs,longtable,array,tabularx}
\usepackage{enumitem}
\usepackage{microtype}
\usepackage{xcolor}
\usepackage{hyperref}
\usepackage{xurl}
\usepackage{fancyhdr}
\usepackage{lastpage}
\usepackage{siunitx}
\usepackage{graphicx}
\usepackage{float}
\usepackage[T1]{fontenc}
\usepackage{lmodern}
\definecolor{navy}{RGB}{25,54,93}
\definecolor{slate}{RGB}{75,85,99}
\hypersetup{colorlinks=true,linkcolor=navy,urlcolor=navy,citecolor=navy,pdftitle={Recommended Reading Order}}
\setlist{nosep,leftmargin=*}
\setlist[description]{style=nextline}
\setlength{\parindent}{0pt}
\setlength{\parskip}{0.52em}
\pagestyle{fancy}
\fancyhf{}
\lhead{\small FMCW / Two-Way Time Transfer Research Corpus}
\rhead{\small Recommended Reading Order}
\cfoot{\small Page \thepage\ of \pageref{LastPage}}
\newcommand{\Saeed}{Saeed}
\newcommand{\AWR}{AWR2944}
\newcommand{\DCA}{DCA1000}
\newcommand{\TWTT}{two-way time transfer}
\newcommand{\FMCW}{FMCW}
\newcommand{\codeword}[1]{\path{#1}}
\newenvironment{keybox}{\begin{center}\begin{minipage}{0.94\textwidth}\hrule\vspace{0.45em}\textbf{Key point.} }{\vspace{0.45em}\hrule\end{minipage}\end{center}}
\begin{document}
\begin{titlepage}
\centering
\vspace*{1.2in}
{\Huge\bfseries\color{navy} Recommended Reading Order\par}
\vspace{0.3in}
{\Large A time-efficient sequence for understanding, implementing, and then extending the FMCW/TWTT system\par}
\vspace{0.6in}
{\large Summer 2026 - Saeed FMCW/TWTT Project\par}
\vfill
{\large Curated research synthesis and engineering reference\par}
{\large Prepared 10 August 2026\par}
\vspace{0.4in}
\begin{minipage}{0.82\textwidth}
\small
This document is a project research aid, not a claim of novelty. It distinguishes source-supported prior art, engineering inference, and proposed simulation steps. Publisher/IEEE PDFs supplied by the user are retained in the archive for private research and are not represented as redistributable.
\end{minipage}
\end{titlepage}
\tableofcontents
\clearpage

\section{If you have 30 minutes}
\begin{enumerate}
\item Read the Saeed meeting notes once, focusing only on the requested progression: ordinary FMCW MATLAB plots first, second radar later, coding after that.
\item Read the first two pages of the 2007 Roehr synchronization paper. Look at the two active stations and how mixing produces equations in time/frequency offset.
\item Read the abstract, Fig. 1, and introduction of the 2019 CFDDS-TWR paper. This shows the mature two-direction version of the idea.
\item Read the first two pages of the Scheiblhofer simulator paper. The block structure is the template for our eventual simulator.
\end{enumerate}

\section{If you have two hours}
Add Scherr 2015 through the FMCW equations/CZT motivation; Ayhan 2016 through the ramp-linearity setup; and the TI chirp programming note. At the end of this session you should be able to derive $f_b=S\tau$, explain why a raw FFT bin is not a precision bound, and map $S,B,T_c,F_s,N$ to radar configuration fields.

\section{Before writing MATLAB}
Read these in order:
\begin{enumerate}
\item Scheiblhofer 2008 simulator - sections on sweep generation, phase integration, mixer/baseband, and noise.
\item Roehr 2007 - synchronization equations using opposite slopes.
\item Scherr 2015 - frequency plus phase evaluation.
\item The Mathematical Handbook in this archive.
\item The MATLAB Simulation Blueprint in this archive.
\end{enumerate}
Then implement only V0/V1.

\section{Before adding nonidealities}
\begin{enumerate}
\item Ayhan 2016: ramp nonlinearity, PN, SNR.
\item Herzel 2018: PLL PN and range correlation.
\item Tschapek 2022: colored phase noise and systematic phase errors.
\item Gottinger 2019: how reciprocal processing handles independent oscillators.
\end{enumerate}

\section{Before touching two physical radars}
Read the AWR2944 datasheet, relevant TRM clock/chirp chapters, EVM guide, TI chirp programming report, calibration note, and DCA1000 capture guide. Make a hardware checklist that distinguishes shared reference frequency from chirp-trigger alignment and RF phase coherence.

\section{Before adding phase coding}
Read Lampel 2020, Uysal 2020, and Kumbul 2022. The question is not only ``what code do we use?'' but ``how do we preserve FMCW's stretch-processing/low-ADC advantage after the code is delayed and misaligned?''

\section{Before discussing novelty}
Read the full 2007 cooperative papers, 2008/2013 cooperative FMCW lineage, 2019 CFDDS-TWR, 2021 loosely coupled bistatic/coherent radar-network work, Honeywell timing/coded FMCW patents, and modern picosecond wireless time-transfer papers. Then formulate a one-sentence claim that contains the exact new mechanism, hardware constraint, estimator, calibration technique, or coded network capability. If the sentence is simply ``FMCW two-way time transfer,'' it is not specific enough.

\section{Suggested personal notes template}
For every important paper, capture six lines:
\begin{enumerate}
\item Problem they solve.
\item Exact waveform/architecture.
\item Unknown parameters estimated.
\item Key equations/estimator.
\item Best demonstrated precision/accuracy and conditions.
\item One limitation or gap that matters to AWR2944/TWTT.
\end{enumerate}
This is a better novelty map than a pile of highlights.

\end{document}
